\clearpage\thispagestyle{empty}
\pdfbookmark[0]{Sources for the May 1805 Juno Report}{juno180505-sources}
\begin{center}\textbf{Sources for the May 1805 Juno Report}\end{center}
\noindent This AI translation follows Franz Xaver von Zach’s report as excerpted in \href{https://gdz.sub.uni-goettingen.de/id/PPN236007467}{Gauss’s \textit{Werke}, vol.~VI} (Göttingen, 1874), pp.~263--265. The report is article XLII, \textit{Fortgesetzte Nachrichten über den neuen Haupt-Planeten Juno}, in the \textit{Monatliche Correspondenz}, vol.~XI, May 1805. The observations by Giuseppe Piazzi at Palermo are on \href{https://www.digitale-sammlungen.de/en/view/bsb10538603?page=489}{p.~475}, followed by the comparison and Carl Friedrich Gauss’s fifth set of elements on \href{https://www.digitale-sammlungen.de/en/view/bsb10538603?page=490}{p.~476}. Gauss’s observation of 20 February and the ephemeris begin on \href{https://www.digitale-sammlungen.de/en/view/bsb10538603?page=491}{p.~477}. The journal also reports on observations from Barnaba Oriani and Wilhelm Olbers and attributes orbital elements to Johann Karl Burckhardt, citing the \textit{Moniteur} of 28 December 1804 on \href{https://www.digitale-sammlungen.de/en/view/bsb10538603?page=496}{p.~482}.
\par\smallskip\noindent Piazzi describes the circle in Book~II of \href{https://www.e-rara.ch/zut/content/titleinfo/11884975}{\textit{Della specola astronomica de’ regi studi di Palermo}} (Palermo, 1792). The Berlin-time elements cited by the collected edition appear on \href{https://www.digitale-sammlungen.de/en/view/bsb10538337?page=280}{p.~270} of Johann Elert Bode’s \href{https://www.digitale-sammlungen.de/en/view/bsb10538337?page=7}{\textit{Astronomisches Jahrbuch für das Jahr 1808}} (Berlin, 1805).
\par\smallskip\noindent Translation and typesetting produced by AI. Journal and yearbook scans from the Bayerische Staatsbibliothek.
\clearpage
